\section{The scope a resolver runs in}
\label{sec:ch03-scopes}
\index{dependency injection!resolver scope|(}

Position thirteen executes the operation, and somewhere inside it your resolver
finally runs. This is the part chapter~\ref{ch:hotchocolate-quickly} got wrong
first time, so it is the part I want to read out of the source rather than
describe.

The counter that chapter needed had to accumulate across a whole request. I
registered it as a scoped service, injected it into the domain services, and got
this:

\begin{minted}{text}
info: Service lookups this request: 1
info: Service lookups this request: 1
info: Service lookups this request: 1
\end{minted}

A hundred and forty-six times. I explained it then by saying HotChocolate
resolves a resolver's services from a scope of its own, one per resolver
invocation. That is true, and it is about sixty per cent of the truth.

\subsection*{Ten lines that explain it}

\index{dependency injection!ResolverTask}%
Here is the code, in \code{ResolverTask.Execute.cs}:

\begin{minted}{csharp}
private async ValueTask ExecuteResolverPipelineAsync(CancellationToken cancellationToken)
{
    if (_context.Field.DependencyInjectionScope == DependencyInjectionScope.Resolver)
    {
        var serviceScope = _operationContext.Services.CreateAsyncScope();
        _context.Services = serviceScope.ServiceProvider;
        _context.RegisterForCleanup(serviceScope.DisposeAsync);
        _operationContext.ServiceScopeInitializer.Initialize(
            _context, _context.RequestServices, _context.Services);
    }

    await _context.ResolverPipeline!(_context).ConfigureAwait(false);
    ...
}
\end{minted}

The scope is created \emph{per field}, and behind an \code{if}. So it is not a
rule about resolvers, it is a property of the field being resolved, and a field
can have the other value. That immediately raises the question of what sets it,
and the answer is the part I did not know:

\begin{minted}{csharp}
public DependencyInjectionScope DefaultQueryDependencyInjectionScope { get; set; } =
    DependencyInjectionScope.Resolver;

public DependencyInjectionScope DefaultMutationDependencyInjectionScope { get; set; } =
    DependencyInjectionScope.Request;
\end{minted}

\index{dependency injection!query and mutation defaults}%
Queries and mutations do not default to the same thing. A query field gets its
own scope per resolver; a mutation field shares the request's. \code{ObjectField}
picks between them when the schema is built, unless the field states its own
preference.

The asymmetry is the reason the setting exists.
Query resolvers run concurrently and mostly independently, and giving each one
its own scope means a scoped \code{DbContext} in one field cannot be dragged
across a thread boundary by another. Mutation resolvers usually want the
opposite: one unit of work, one transaction, one \code{DbContext} that several
steps commit together. Chapter~\ref{ch:data-without-the-n-plus-1} is where this
stops being trivia, because a \code{DbContext} is exactly the service whose
scope you have to be able to reason about.

\subsection*{Two attributes, and a way to see them work}

\index{dependency injection!UseRequestScope}%
Any field can override the default. In the descriptor API those are
\code{UseRequestScope()} and \code{UseResolverScope()}; Mosaic is
implementation-first, so the relevant form is the pair of attributes,
\code{[UseRequestScope]} and \code{[UseResolverScope]}, both of which go on a
method.

The companion repository has a sample that does nothing but show the difference.
A scoped service hands out a number that increments every time a scope
constructs one, and two fields report both the number they were given and the
number living in the request's own provider:

\begin{minted}{csharp}
[QueryType]
public static partial class ScopeQueries
{
    public static ScopeReading GetDefaultScope(
        ScopeProbe injected,
        IResolverContext context)
        => new(
            injected.Id,
            context.RequestServices.GetRequiredService<ScopeProbe>().Id);

    [UseRequestScope]
    public static ScopeReading GetRequestScope(
        ScopeProbe injected,
        IResolverContext context)
        => new(
            injected.Id,
            context.RequestServices.GetRequiredService<ScopeProbe>().Id);
}
\end{minted}

Ask for both, twice each, in a single request:

\begin{minted}{graphql}
{
  a: defaultScope { injected fromRequestServices }
  b: defaultScope { injected fromRequestServices }
  c: requestScope { injected fromRequestServices }
  d: requestScope { injected fromRequestServices }
}
\end{minted}

\begin{minted}[fontsize=\footnotesize]{json}
{"data":{
  "a":{"injected":1,"fromRequestServices":2},
  "b":{"injected":3,"fromRequestServices":2},
  "c":{"injected":2,"fromRequestServices":2},
  "d":{"injected":2,"fromRequestServices":2}}}
\end{minted}

Read the numbers rather than the field names. There are three probes in that
request, not one. Field \code{a} was handed probe~1 and field \code{b} probe~3,
and neither is probe~2, which is the one the request itself owns. The two
annotated fields were both handed probe~2. Send the query again and the same
pattern comes back with 4, 5 and 6.

The ordering is a small bonus: 1 before 2 before 3 says the first resolver's
private scope was created before anything asked the request for its copy.

\begin{figure}[htbp]
  \centering
  % Request scope versus the child scopes a query field gets by default, with the
% probe numbers the sample actually returned. Referenced from chapter 03. Bare
% tikzpicture: the figure environment, caption and label live at the call site.
% Uses only arrows.meta and positioning.
\begin{tikzpicture}[
    font=\small,
    outer/.style={draw, thick, rounded corners=3pt, minimum width=104mm,
                  minimum height=46mm},
    shared/.style={draw, thick, rounded corners=2pt, minimum width=44mm,
                   minimum height=11mm, align=center, font=\scriptsize},
    child/.style={draw, rounded corners=2pt, minimum width=30mm,
                  minimum height=15mm, align=center, font=\scriptsize},
    plain/.style={draw, dashed, rounded corners=2pt, minimum width=30mm,
                  minimum height=15mm, align=center, font=\scriptsize},
    reach/.style={-{Stealth[length=2mm]}, dashed},
    note/.style={font=\scriptsize\itshape, align=center}
  ]

  \node[outer] (req) at (0,0) {};

  \node[note, anchor=south] at (0,2.55)
    {\textbf{the request scope}, which is what \texttt{RequestServices} points at};

  \node[shared] (rprobe) at (0,1.25)
    {the request's own probe:\\\textbf{2}};

  \node[child] (fa) at (-3.4,-1.1)
    {field \texttt{a}\\its own child scope\\was handed \textbf{1}};

  \node[child] (fb) at (0,-1.1)
    {field \texttt{b}\\its own child scope\\was handed \textbf{3}};

  \node[plain] (fc) at (3.4,-1.1)
    {field \texttt{c}, marked\\\texttt{[UseRequestScope]}\\was handed \textbf{2}};

  \draw[reach] (fc) -- (rprobe);

  \node[note, anchor=north] at (0,-2.75)
    {Three probes, one request. Only the annotated field reaches the one the
     request owns.};

\end{tikzpicture}

  \caption{Why the counter reported one. Each query field resolved by default
  gets a child scope, so a scoped service injected into it is constructed fresh
  and disposed when the field finishes. \code{RequestServices} always points at
  the request's provider, and \code{[UseRequestScope]} makes an injected
  parameter come from there too.}
  \label{fig:ch03-scopes}
\end{figure}

\subsection*{What that means for the counter}

\index{dependency injection!RequestServices}%
So chapter~\ref{ch:hotchocolate-quickly} had three fixes available and I picked
one without knowing the others existed.

Hanging the total off \code{HttpContext.Items}, which is what Mosaic does,
works and has the advantage of being obvious to anyone who knows ASP.NET Core
and nothing about HotChocolate. Annotating every contributing field with
\code{[UseRequestScope]} would also have worked, and would have been worse: it
is an attribute on every resolver that touches a domain service, and one that
has to be remembered on the next one somebody adds. The third option is the one
this chapter has been leaning on the whole time, and it is what the timeline
uses.

\code{IResolverContext.RequestServices} is not the resolver's scope. Look again
at the ten lines above. The \code{if} replaces the resolver context's
\code{Services} property and never touches the other one, which comes from the
operation context. That was initialised from the request, so it is the
request's provider whatever the field's scope setting says. A service
registered scoped and read through it is therefore one instance per request,
reachable from any resolver, with no attributes and no \code{HttpContext}:

\begin{minted}[fontsize=\footnotesize]{csharp}
public override IDisposable ResolveFieldValue(IMiddlewareContext context)
{
    context.RequestServices.GetService<RequestTimeline>()?.CountResolver();
    return EmptyScope;
}
\end{minted}

That is how a hundred and forty-six resolver invocations, many of them
concurrent, all increment the same counter. Which leaves the question of where
that method comes from.
\index{dependency injection!resolver scope|)}
